  % ======================================================================
  %  DataGatherer --- Technical assessment for industry readers
  %  Based on the personal template in docs/template/doc.tex
  %
  %  Compile: latexmk -pdf industrial_report.tex
  % ======================================================================

  \documentclass[10pt,a4paper]{scrartcl}

  % --------------------------------------------------
  % Encoding and language
  % --------------------------------------------------
  \usepackage[T1]{fontenc}
  \usepackage[utf8]{inputenc}
  \usepackage[english]{babel}

  % --------------------------------------------------
  % Page layout
  % --------------------------------------------------
  \usepackage[
    top=18mm,
    bottom=18mm,
    left=20mm,
    right=20mm,
    includefoot
  ]{geometry}

  % --------------------------------------------------
  % Mathematics and graphics
  % --------------------------------------------------
  \usepackage{graphicx}
  \usepackage{amsmath,amssymb}
  \usepackage{bm}
  \setlength{\jot}{0.5cm}

  % --------------------------------------------------
  % Tables
  % --------------------------------------------------
  \usepackage{booktabs}
  \usepackage{tabularx}
  \usepackage{array}
  \usepackage{ragged2e}

  % A ragged-right fixed-width column for use in tabularx:
  %   \begin{tabularx}{\textwidth}{@{}L{3cm} X@{}}
  \newcolumntype{L}[1]{>{\RaggedRight\arraybackslash}p{#1}}

  % --------------------------------------------------
  % Typography
  % --------------------------------------------------
  \usepackage[scaled=0.95]{helvet}
  \usepackage[expansion=false]{microtype}   % expansion=false avoids font errors
  \renewcommand{\familydefault}{\sfdefault} % sans-serif throughout

  % --------------------------------------------------
  % Lists
  % --------------------------------------------------
  \usepackage{enumitem}
  \setlist[itemize]{
    leftmargin=*,
    topsep=2pt,
    parsep=0pt,
    partopsep=0pt,
    itemsep=3pt
  }
  \setlist[description]{
    leftmargin=0pt,
    topsep=2pt,
    parsep=0pt,
    itemsep=3pt
  }

  % --------------------------------------------------
  % Colour and hyperlinks
  % --------------------------------------------------
  \usepackage[dvipsnames]{xcolor}
  \usepackage[hidelinks]{hyperref}

  \definecolor{accent}{RGB}{0,129,88}   % change this to restyle everything

  % \emph produces upright accent-coloured text instead of italic.
  % Use \textit{} when you genuinely want italics.
  \renewcommand{\emph}[1]{\textcolor{accent}{#1}}

  % --------------------------------------------------
  % Header and footer
  % --------------------------------------------------
  \usepackage[automark,headsepline=false,footsepline=false]{scrlayer-scrpage}
  \clearpairofpagestyles

  \ifoot{Victor Gabriele De Matteis}
  \cfoot{\pagemark}
  \ofoot{DataGatherer --- Technical assessment}

  \setlength{\footskip}{20mm}  % pushes footer toward the page edge
  \pagestyle{scrheadings}

  % --------------------------------------------------
  % Spacing
  % --------------------------------------------------
  \setlength{\parindent}{0pt}
  \setlength{\parskip}{0.45em}

  % --------------------------------------------------
  % Section formatting
  % --------------------------------------------------
  \RedeclareSectionCommand[
    beforeskip=1.4ex plus 0.2ex minus 0.1ex,
    afterskip=0.6ex
  ]{section}
  \RedeclareSectionCommand[
    beforeskip=1.0ex plus 0.2ex minus 0.1ex,
    afterskip=0.4ex
  ]{subsection}
  \RedeclareSectionCommand[
    beforeskip=0.7ex plus 0.2ex minus 0.1ex,
    afterskip=0.3ex
  ]{subsubsection}

  \setkomafont{section}{\large\bfseries\color{accent}}
  \setkomafont{subsection}{\normalsize\bfseries\color{accent}}
  \setkomafont{subsubsection}{\normalsize\bfseries}

  % A section command that adds a faint rule beneath the title.
  % Use \RuleSection{Title} instead of \section{Title} where you want it.
  \newcommand{\RuleSection}[1]{%
    \section{#1}%
    \vspace{0.2em}{\color{accent!35}\hrule height 0.4pt}\vspace{0.3em}%
  }

  % --------------------------------------------------
  % Additions for this document
  % --------------------------------------------------
  \DisableLigatures[-]{encoding=T1, family=tt*}
  \DeclareRobustCommand{\code}[1]{{\frenchspacing\texttt{\detokenize{#1}}}}

  % Readiness ratings
  \newcommand{\ready}{\textcolor{accent}{\textbf{Ready}}}
  \newcommand{\partly}{\textcolor{Dandelion!80!black}{\textbf{Partial}}}
  \newcommand{\notready}{\textcolor{BrickRed}{\textbf{Not ready}}}

  \setlength{\emergencystretch}{3em}
  \usepackage[section]{placeins}
  \makeatletter
  \setlength{\@fptop}{0pt}
  \makeatother

  % ======================================================================
  \begin{document}
  % ======================================================================

  % --------------------------------------------------
  % Title block
  % --------------------------------------------------
  \begin{center}
    {\Huge\bfseries\color{accent} DataGatherer}\\[0.25cm]
    {\small Technical assessment --- status, risks and next steps --- October 2026}\\[0.15cm]
    {\small Victor Gabriele De Matteis \quad\textbullet\quad github.com/ViDem02/DataGatherer}\\[0.2cm]
    {\color{accent!35}\hrule height 0.5pt}
  \end{center}

  \vspace{0.3em}

  % ======================================================================
  %  BODY
  % ======================================================================

  \RuleSection{Summary}
  \label{sec:summary}

  \textbf{What it is.} A web application that turns screenshots of Instagram posts into a spreadsheet of who
  used which hashtag, how many times. An operator crops the account name and the hashtags on each screenshot,
  text recognition (OCR) reads them, the operator confirms or corrects the text, and the application builds
  a user--hashtag matrix ready for network or market analysis.

  \textbf{Status.} Suitable \emph{today} as a tool run on one computer by one person or a small team. It is
  \emph{not yet suitable} to be offered to customers or hosted on a network: the public code history still
  holds old secrets and third-party personal data, the collected data has no GDPR framework, and there is no
  HTTPS, automated backup or automated test gate.

  \textbf{Recommendation.} Use it internally as is. Before any wider use, complete the four items of
  Section~\ref{sec:next} marked ``before any external use''; together they are estimated at one to two weeks.

  \begin{table}[h]
  \small
  \centering
  \renewcommand{\arraystretch}{1.1}
  \begin{tabularx}{\textwidth}{@{}L{4.4cm} L{2cm} X@{}}
  \toprule
  \textbf{Area} & \textbf{Rating} & \textbf{Basis} \\[.3cm]
  \midrule
  Core functionality & \ready & Upload, cropping, OCR, verification, matrix export all work end to end. \\[.3cm]
  Ease of installation & \ready & Docker Desktop plus a double-click; ready-made images (about 780\,MB; 36 seconds to download in our test). \\[.3cm]
  Application security & \ready & Reviewed and hardened in October 2026: per-user data isolation, secure login, login lockout, private files. \\[.3cm]
  Repository hygiene & \notready & Old secrets and personal data remain in the public git history. \\[.3cm]
  Data protection (GDPR) & \notready & Personal data of third parties is collected without legal basis, retention or deletion process. \\[.3cm]
  Reliability and operations & \partly & No automated backups; long batch computations stop if the server restarts. \\[.3cm]
  Quality assurance & \partly & 39 automated tests, run by hand; CI builds the images but does not run the tests yet. \\[.3cm]
  Scalability & \partly & One server; fine for thousands of images, not for many simultaneous teams. \\[.3cm]
  \bottomrule
  \end{tabularx}
  \caption{Readiness at a glance.}
  \label{tab:readiness}
  \end{table}

  \RuleSection{Value}
  \label{sec:value}

  The repository's own measurements (\code{stats.md}) report 1{,}089 screenshots processed. They estimate about
  2 seconds of automated work per image against about 30 seconds for manual transcription, i.e.\ roughly 93\,\%
  less effort. These figures are \emph{estimates}, not timed measurements, and they exclude the operator's
  verification of each text, which remains necessary because OCR misreads accents and similar letters.

  The output is a standard CSV file (rows: hashtags, columns: accounts, cells: number of posts), directly usable in
  spreadsheets, Python/R or network-analysis tools.

  \RuleSection{How It Is Built}
  \label{sec:built}

  \begin{table}[h]
  \small
  \centering
  \renewcommand{\arraystretch}{1.1}
  \begin{tabularx}{\textwidth}{@{}L{3.4cm} X@{}}
  \toprule
  \textbf{Component} & \textbf{Technology} \\[.3cm]
  \midrule
  User interface & React web app, served by nginx; works in any modern browser. \\[.3cm]
  Application & Python, Django 5.2 LTS (supported until April 2028), REST API, Tesseract OCR. \\[.3cm]
  Data & PostgreSQL 16; uploaded images on a Docker volume. \\[.3cm]
  Delivery & Three Docker containers; images built automatically for Intel/AMD and ARM and published on GitHub. \\[.3cm]
  Runs on & Windows, macOS or Linux with Docker Desktop. \\[.3cm]
  \bottomrule
  \end{tabularx}
  \caption{Technology.}
  \label{tab:tech}
  \end{table}

  The application is reachable only from the computer it runs on, by design. Each user account sees only its own
  projects and images; an administrator manages accounts through a built-in admin panel.

  \RuleSection{What Was Done in October 2026}
  \label{sec:done}

  A review on 12 October 2026 found the application unsafe to expose: anyone could download the whole database,
  users could read and change each other's data, and secrets were stored in the code. The same day:
  \begin{itemize}
    \item \textbf{Security:} every hole found was closed; login tokens moved to browser-protected cookies; repeated
          wrong passwords lock the account for 15 minutes; uploaded images became private; dependency scans now report
          no known vulnerabilities (they reported 64 in the backend and 116 in the frontend).
    \item \textbf{Installation:} from a multi-step technical setup (Node.js, secrets written by hand, 10--20 minutes of building) to
          a double-click that generates its own secrets and downloads ready-made images.
    \item \textbf{Correctness:} batch jobs no longer report success when part of the work failed; the matrix export
          went from one database query per cell to three queries in total (a 20 $\times$ 20 example: 442 queries
          before); edits that fail can no longer delete existing data.
  \end{itemize}
  Each fix was verified with automated tests or end-to-end checks, including a real browser.

  \RuleSection{Risks}
  \label{sec:risks}

  \begin{table}[h]
  \small
  \centering
  \renewcommand{\arraystretch}{1.1}
  \begin{tabularx}{\textwidth}{@{}L{4.2cm} L{1.5cm} X@{}}
  \toprule
  \textbf{Risk} & \textbf{Level} & \textbf{Mitigation} \\[.3cm]
  \midrule
  Old secrets, database dumps and a third party's utility bill in the public code history & High & Rotate every secret (done locally), then rewrite the history or recreate the repository. \\[.3cm]
  Personal data of Instagram users (GDPR); platform terms of use & High & Legal review: legal basis, retention period, deletion on request, record of processing; check Instagram's terms. \\[.3cm]
  Data loss & Medium & Scheduled database and image backups with a tested restore. \\[.3cm]
  Wrong numbers in the matrix from operator errors (duplicate screenshots, OCR typos accepted) & Medium & Guards in the interface: duplicate detection, warnings before saving (designed, not built). \\[.3cm]
  Regressions from future changes & Medium & Run the tests automatically on every change. \\[.3cm]
  Exposure if run on a network & Medium & Keep it local, or add HTTPS and secure cookies before exposing it. \\[.3cm]
  Windows installation not yet tried on a real PC & Low & One test run on Windows. \\[.3cm]
  \bottomrule
  \end{tabularx}
  \caption{Main risks.}
  \label{tab:risks}
  \end{table}

  \clearpage
  \RuleSection{Next Steps}
  \label{sec:next}

  \begin{table}[h]
  \small
  \centering
  \renewcommand{\arraystretch}{1.1}
  \begin{tabularx}{\textwidth}{@{}L{6.4cm} L{2.4cm} X@{}}
  \toprule
  \textbf{Action} & \textbf{Effort} & \textbf{When} \\[.3cm]
  \midrule
  Purge the code history and rotate any remaining secret & 1 day & Before any external use \\[.3cm]
  GDPR assessment of the collected data & legal, 1--2 weeks & Before any external use \\[.3cm]
  Automated tests on every change, scan for committed secrets & 2--3 days & Before any external use \\[.3cm]
  Automated backups and a tested restore & 1--2 days & Before any external use \\[.3cm]
  Operator guards against duplicate screenshots and OCR typos & 2--3 weeks & To improve data quality \\[.3cm]
  HTTPS, background job queue, monitoring, staging environment & 2--4 weeks & Only for a hosted, multi-team service \\[.3cm]
  \bottomrule
  \end{tabularx}
  \caption{Recommended actions. Efforts are estimates for one developer.}
  \label{tab:next}
  \end{table}

  Detailed technical documentation is available on request: a change record of the security work, the user
  and maintenance workflows, and an in-depth description of the architecture and algorithms.

  \end{document}
